\textbf{Video Retrieval.}
We evaluate \model on the video retrieval task. Given a set of videos and related natural language captions, this task requires retrieving the matched video or caption corresponding to its inter-modality counterpart from candidates. We follow the common paradigm to capture visual and text semantics by a visual encoder $f_v(\cdot)$ and a text encoder $f_t(\cdot)$, then calculate the cross-modality similarity matrices as the retrieval guidance. We leverage the multimodal video encoder as $f_v(\cdot)$ and $f_t(\cdot)$ with pretrained ViT-L/14~\cite{dosovitskiy2020image} as the basic CLIP~\cite{CLIP} architecture and finetune the entire model on each retrieval dataset. The training recipes and most of the hyperparameter settings follow CLIP4Clip~\cite{clip4clip}, including training schedule, learning rate, batch size, video frames, maximum text length, etc. To boost model performance, we also adopt the dual softmax loss~\cite{cheng2021improving} as the post-processing operation.

Our model is evaluated on six public benchmarks: MSR-VTT~\cite{xu2016msr}, MSVD~\cite{wu2017deep}, LSMDC~\cite{anne2017localizing}, DiDeMo~\cite{rohrbach2015dataset}, ActivityNet~\cite{caba2015activitynet}, and VATEX~\cite{wang2019vatex}, where we report the results on the standard split following previous works. We measure the retrieval results under the rank-1 (R@1) metric both on text-to-video and video-to-text tasks, which are shown in Table~\ref{tab:video_retrieval}. Results show that our model significantly outperforms all previous methods by a large margin, showing the superiority of \model~ on video-language related tasks. More detailed retrieval results, including rank-5 (R@5) and rank-10 (R@10), can be found in the supplementary materials.
